\documentclass[conference]{IEEEtran}
\usepackage{cite}
\usepackage{amsmath,amssymb,amsfonts}
\usepackage{graphicx}
\usepackage{textcomp}
\usepackage{xcolor}
\usepackage{booktabs}
\usepackage{hyperref}
\usepackage{array}

\begin{document}

\title{WaterWallet: A Physics-Informed Decision-Support Platform for\\
Residential Rainwater Harvesting in the Mumbai Metropolitan Region}

\author{
  \IEEEauthorblockN{Nidhish Shirur, Neha Bhatia, Diksha Bagwe}
  \IEEEauthorblockA{
    \textit{Department of Computer Engineering}\\
    \textit{Vivekanand Education Society's Institute of Technology (VESIT)}\\
    University of Mumbai, Mumbai, India\\
    Guide: Dr.~Mrs.~Gresha Bhatia
  }
}

\maketitle

\begin{abstract}
Residential housing societies need to assess whether rooftop rainwater harvesting (RWH) can contribute to non-potable water demand and whether a proposed system is financially plausible. This paper presents WaterWallet, a full-stack decision-support prototype that combines society demand and billing inputs, map-derived rooftop catchments, historical rainfall retrieval, seasonal rainfall forecasting, physics-based yield calculations, hydrogeological screening, and financial analysis. Its React and Vite client communicates with a FastAPI service layer. The backend includes a deterministic RWH safety gate, bundled Random Forest models for configuration recommendation and monthly yield estimation, a daily tank water-balance sizing service, deterministic and Monte Carlo return-on-investment (ROI) analyses, and an optional natural-language advisory. The system is designed for pre-consultation feasibility analysis rather than final engineering design. Its classifier is trained on rule-generated synthetic labels, the yield regressor learns targets generated from the physical yield equation, and the rainfall forecast has not been independently validated here against a current holdout evaluation. Accordingly, no model-accuracy, field-savings, or generalization claim is made. Existing software tests exercise selected water-balance and scenario-flow behaviors; they do not establish hydrological validity or field performance. The paper describes the implemented methods and their evidence limits to provide a reproducible, technically bounded account of the current prototype.
\end{abstract}

\begin{IEEEkeywords}
Rainwater harvesting, decision-support system, urban hydrology,
seasonal forecasting, water balance, Monte Carlo simulation,
explainable machine learning, Mumbai Metropolitan Region.
\end{IEEEkeywords}

\section{Introduction}
\label{sec:intro}

Rainwater harvesting can supplement municipal and tanker water for
non-potable uses in residential buildings. However, a society committee
considering an installation must reason about more than annual rainfall:
roof area and material affect runoff, demand varies over time, storage
limits determine how much monsoon inflow can be retained, and site
conditions constrain groundwater recharge. Financial estimates are also
sensitive to system cost, water tariffs, tanker volume, and maintenance.

WaterWallet addresses this decision problem as a feasibility-screening
workflow. It combines a society demand profile and billing history with a
rooftop polygon, rainfall information, site attributes, and configurable
financial assumptions. The interface presents a catchment estimate,
configuration recommendation, tank-sizing analysis, and ROI projections.
It does not certify that a site is safe for recharge, guarantee a forecast,
or replace a licensed engineer's site investigation and final design.

The current implementation makes three practical contributions:
\begin{enumerate}
  \item It connects an interactive frontend to a backend API that performs
    demand, rainfall, yield, recommendation, tank-sizing, and ROI
    calculations in one user workflow.
  \item It separates deterministic physical and safety rules from
    data-driven recommendations, and distinguishes annualized financial
    estimates from storage-constrained water-balance simulations.
  \item It documents the evidentiary limits of its predictive components:
    the configuration labels are synthetic and rule-derived, the yield
    model's training targets come from a physical equation, and the
    available software tests are not field validation.
\end{enumerate}

\section{Related Methods and Design Context}
\label{sec:related}

WaterWallet uses the 135 litres per capita per day residential baseline
configured from IS~1172 \cite{bis1172}. This value is an input convention
for the prototype, not a measurement of any particular society's
consumption. The frontend and backend also permit an optional
unaccounted-for-water (UFW) allowance and a peak-demand buffer; these
adjustments are application settings and should not be interpreted as a
site-specific demand study.

For rainfall, the application queries the Open-Meteo Archive API
\cite{openmeteo}. Its forecasting endpoint applies seasonal Auto-ARIMA
model selection to a chronological monthly series. Seasonal time-series
methods are established forecasting tools \cite{hyndman}; their inclusion
does not by itself establish forecast skill for a particular location or
horizon. The service returns a forecast interval, but a current
independent evaluation is outside the evidence reported in this paper.

The recommendation service can return per-feature SHAP attributions
\cite{lundberg}. Such attributions explain the model output with respect
to the fitted model; they are neither causal effects nor evidence that a
recommendation is hydrologically correct. WaterWallet therefore places a
deterministic safety gate before the learned classifier and treats both
the gate and the model as screening aids that require expert review.

\section{System Architecture}
\label{sec:architecture}

\begin{figure*}[t]
\centering
\fbox{\begin{minipage}{0.94\textwidth}
\centering
\textbf{React 19 + Vite client}\\
Virtual Meter: society demand and billing \qquad
Map Studio: rooftop polygon and location \qquad
ROI Dashboard: recommendation, tank sizing, and financial views\\[4pt]
$\Downarrow$ Axios JSON requests; shared workflow state in React Context\\[4pt]
\textbf{FastAPI application}\\
Demand and society routes \quad Catchment and forecast routes \quad
ML recommendation and yield routes \quad ROI and tank-sizing routes\\[4pt]
$\Downarrow$\\[4pt]
\textbf{Calculation and data services}\\
Demand and yield equations \quad Open-Meteo rainfall \quad
Safety gate and bundled ML artifacts \quad Daily and monthly water balances\\
Deterministic ROI and Monte Carlo scenarios\\[4pt]
\textbf{Optional integrations:} Esri World Imagery tiles and Groq advisory\\
\textbf{Persistence boundary:} PostgreSQL schema and selected database paths;
ordinary society CRUD remains in an in-memory store
\end{minipage}}
\caption{Implemented WaterWallet request and service architecture.
Optional integrations and persistence paths are distinguished from the
core calculation flow.}
\label{fig:architecture}
\end{figure*}

\subsection{Client and Service Organization}

The client is implemented with React and Vite. Its principal views are
the Virtual Meter, Map Studio, and ROI Dashboard. React Context carries
society, catchment, rainfall, demand, and tank-capacity values between
views. The map view uses Leaflet and Esri World Imagery tiles; Turf.js
calculates polygon area from the traced boundary. API requests are made
through an Axios client to the FastAPI backend.

The backend registers routers for authentication, society and demand
operations, geospatial catchment calculations, ML inference, ROI and tank
sizing, and advisory generation. Pydantic schemas validate request and
response data. The current calculation services and their principal
interfaces are summarized in Table~\ref{tab:api}.

\begin{table}[htbp]
\caption{Principal Backend Interfaces}
\label{tab:api}
\centering
\small
\begin{tabular}{p{0.39\columnwidth}p{0.53\columnwidth}}
\toprule
\textbf{Endpoint} & \textbf{Role in the workflow}\\
\midrule
\url{/api/auth/*} & Register, authenticate, and retrieve the current user\\
\url{/api/societies/baseline} & Calculate total and potable/non-potable demand\\
\url{/api/catchment/calculate} & Retrieve monthly rainfall and calculate yield\\
\url{/api/catchment/forecast} & Forecast rainfall and forecast-based yield\\
\url{/api/ml/recommend} & Apply the safety gate, then classify configuration\\
\url{/api/ml/yield} & Estimate monthly yield with physics fallback\\
\url{/api/roi/simulate} & Calculate deterministic multi-year ROI\\
\url{/api/roi/monte-carlo} & Estimate a break-even distribution\\
\url{/api/roi/tank-sizer} & Compare candidate tanks using daily rainfall\\
\url{/api/advisory/water} & Request an optional Groq-generated advisory\\
\bottomrule
\end{tabular}
\end{table}

\subsection{Persistence and External Dependencies}

SQLAlchemy models and an Alembic migration define PostgreSQL tables for
users, societies, billing entries, and analysis snapshots. Registration
and login use the database-backed authentication path; passwords are
hashed with bcrypt and access tokens use JWT. The frontend also supports
a guest path. This database support is not equivalent to end-to-end
persistent society management: ordinary society CRUD and billing routes
currently use an authenticated, in-memory store, while explicitly named
database endpoints provide selected read and analysis operations.

Open-Meteo supplies rainfall data, and Esri supplies map tiles. The
advisory route is optional, requires a configured Groq API key, and sends
the submitted society metrics to that external service. The advisory is
prompt-constrained to return structured text; it is not a validated
engineering recommendation. Core deterministic calculations do not
require the advisory service.

\section{Methods}
\label{sec:methods}

\subsection{Demand Estimation}
\label{sec:demand}

Let $F$ denote the number of flats and $O$ the average occupants per flat.
The implementation converts their product to an integer occupant count
$N$ and uses the configured IS~1172 baseline $q=135$ L/person/day:
\begin{align}
  N &= \lfloor F O \rfloor,\\
  D_0 &= N q,\\
  D &= \operatorname{round}\!\left[
      D_0(1 + 0.15 I_{\mathrm{UFW}})(1 + 0.10)
      \right],
  \label{eq:demand}
\end{align}
where $I_{\mathrm{UFW}}$ is 1 when the optional 15\% UFW allowance is
enabled and 0 otherwise. The 10\% demand buffer is applied in either
case. The adjusted demand is divided into potable and non-potable
components using the configured 70:65 LPCD proportions; the non-potable
component is the intended RWH demand target. Annual demand is calculated
as the daily value multiplied by 365. For monthly views, the frontend
distributes annual demand using configured seasonal weights and
normalizes the weights so their sum equals the annual total.

The Virtual Meter accepts society and billing inputs, including
municipal bills, tanker counts, tanker unit costs, and tanker capacities.
These values provide financial context; the billing interface does not
constitute an independently audited consumption record.

\subsection{Catchment, Rainfall, and Yield}
\label{sec:yield}

The user traces a rooftop polygon on the map. Turf.js computes its
geodesic area, and the selected location is used for rainfall retrieval.
For monthly rainfall depth $P_m$ (mm), gross area $A_g$ (m$^2$), roof
runoff coefficient $C_r$, and usable-area factor $\eta=0.80$, the
physics-based estimate is
\begin{equation}
  Y_m = \frac{P_m A_g C_r \eta}{1000}
  \quad \mathrm{kL},
  \label{eq:yield}
\end{equation}
using the equivalence of 1 mm of rainfall on 1 m$^2$ to 1 litre. The
configured runoff coefficients are 0.90 for concrete and metal, 0.75 for
tiles, and 0.45 for gravel. The area factor is intended to account for
roof equipment and other unusable surface; these coefficients are
assumptions rather than measurements of runoff at the user's building.

The catchment endpoint averages monthly totals across the configured
rolling set of 15 complete calendar years. If that monthly retrieval
fails, the service logs the failure and uses a regional monthly
climatological fallback. The daily rainfall endpoint used for tank sizing
has different failure behavior: it rejects malformed or incomplete
data, omits incomplete calendar years, and reports an error when no
complete year is available. These two data paths should not be conflated.

\subsection{Seasonal Rainfall Forecast}
\label{sec:forecast}

The forecast endpoint preserves chronological monthly totals for the
same configured historical period and requires at least 36 observations.
It fits a seasonal Auto-ARIMA model with period 12, selecting orders
through an AIC-based search, and returns between 1 and 24 monthly
forecasts with confidence intervals and the selected model orders.
Forecast-derived yield is calculated with Eq.~(\ref{eq:yield}).

If historical monthly retrieval fails, the service constructs a
deterministic regional synthetic history from fallback monthly normals
and fits the forecasting model to that history. The forecast response
exposes the history source and a fallback flag; the interface labels a
forecast based on synthetic regional history. This provenance improves
interpretability but is not evidence of forecast accuracy.

\subsection{Configuration Recommendation and Safety Gate}
\label{sec:classifier}

The recommendation endpoint offers three configuration labels:
\textit{Storage-Only}, \textit{Borewell-Recharge}, and \textit{Hybrid}.
Before loading the classifier, the deterministic Tier~1 gate forces a
Storage-Only result when any configured check is triggered: coastal
distance is at most 2.5 km; no borewell is present; soil permeability
class is below 2; or a present borewell is shallower than 30 m. A gate
result includes a reason and does not invoke ML or SHAP. These thresholds
are implemented screening rules; this paper does not assert that they
are universally applicable hydrogeological standards.

When the gate passes, Tier~2 uses a bundled calibrated Random Forest
classifier. Its seven inputs are roof area, borewell presence and depth,
soil permeability class, coastal distance, annual rainfall, and peak
daily rainfall. Training code generates 600 synthetic site profiles and
assigns each label through a deterministic rule based on site
attributes. The classifier therefore approximates those labeling rules;
it has not been trained on a representative dataset of independently
surveyed housing societies. SHAP contributions are returned when the
explanation dependency and model are available. The response carries an
explicit warning when explanations are unavailable, and feature
attributions are not presented as causal explanations.

\subsection{Monthly Yield Regressor}
\label{sec:yield-model}

The backend also loads a Random Forest regressor for monthly yield. The
training script constructs inputs from historical monthly rainfall and
multiple roof areas, then generates its target from the same
runoff-and-usability equation in Eq.~(\ref{eq:yield}). At inference, the
API supplies twelve monthly rainfall values and corresponding lag,
seasonal, and area features. If model files or their compatible
scikit-learn version are unavailable, the service falls back to the
physics equation and returns the selected source in its response.

Because the training target is equation-derived rather than measured
catchment yield, the regressor is a learned approximation to the
configured physics calculation. Its presence does not demonstrate
improved prediction over the equation, and no predictive-performance
metric is claimed here.

\subsection{Daily Tank Sizing}
\label{sec:tank}

Tank sizing uses daily rainfall grouped by complete calendar year. For
each candidate capacity $K$, the service simulates each day in sequence.
Daily inflow is
\begin{equation}
  I_d = \max(0,P_d) A_u C_r (1-f_{\mathrm{ff}}),
  \label{eq:inflow}
\end{equation}
where $P_d$ is daily rainfall, $A_u$ is usable roof area, and
$f_{\mathrm{ff}}$ is the first-flush fraction. The tank serves eligible
non-potable demand $D_d$ from current storage plus inflow:
\begin{align}
  S_d &= \min(K_d + I_d,D_d),\\
  K_{d+1} &= \min\!\left(K,\max(0,K_d+I_d-D_d)\right),
\end{align}
where $K_d$ is storage at the start of the day and $S_d$ is useful
rainwater served. Overflow is tracked separately. When recharge is
enabled, overflow up to the configured daily infiltration capacity is
recorded as recharge; excess is recorded as lost overflow. Neither
recharge nor lost overflow is counted as financial savings in this
service.

The endpoint compares caller-supplied capacities (default candidates
range from 5,000 to 70,000 L). It selects the smallest candidate meeting
the requested average reliability when possible. If not, it applies a
fallback reliability target or a diminishing-returns capacity heuristic.
The default 40\% reliability target, costs, recharge capacity, and
diminishing-returns threshold are prototype parameters, not universal
design standards. The service reports the candidate metrics, selection
status, and years used so a user can inspect the basis of the result.

\subsection{Financial Analysis}
\label{sec:roi}

The deterministic ROI service computes an annualized upper-bound
estimate: it caps annual calculated yield at annual demand but does not
simulate within-year tank storage. Eligible yield offsets tanker volume
first and then municipal volume. Gross savings are valued at the current
tanker and municipal rates, with tanker escalation supplied by the caller
and municipal escalation set to 5\% per year. Annual maintenance is
modeled as 3\% of installation cost:
\begin{align}
  G_y &= V_{\mathrm{tank},y}R_{\mathrm{tank},y}
       + V_{\mathrm{mun},y}R_{\mathrm{mun},y},\\
  N_y &= G_y - 0.03\,C_{\mathrm{install}},\\
  B_y &= \sum_{j=1}^{y} N_j - C_{\mathrm{install}},
  \label{eq:roi}
\end{align}
where $G_y$ is gross annual savings, $N_y$ is net annual savings, and
$B_y$ is cumulative net benefit. The service reports the first
break-even point it finds, yearly savings, and cumulative expenditure
paths. Its assumptions make it distinct from the storage-constrained
simulations below.

The Monte Carlo service samples 1,000 scenarios by default. If a
site-specific monthly rainfall baseline is supplied, each month's value
is sampled independently and uniformly from 80\% to 120\% of its
baseline; otherwise the service uses a saved monthly percentile table
when present or a regional fallback distribution. The current service
uses a fixed pseudorandom seed. Monthly inflow is converted to yield
using roof area, roof-surface runoff, and usable-area factors. A
month-by-month tank balance carries storage within the sampled year and
caps it at tank capacity. The resulting annual useful volume is valued
over a ten-year financial horizon with tanker escalation, municipal
escalation, and maintenance costs. The response contains break-even
probabilities, percentiles, never-break-even share, and scenario-source
indicators. Monthly values are sampled independently, and this engine
does not model inter-month statistical dependence, multi-year drought
sequences, or recharge-derived savings.

\section{Implementation and Operational Requirements}
\label{sec:implementation}

Table~\ref{tab:stack} summarizes the principal implementation
technologies and their roles.

\begin{table}[htbp]
\caption{Implementation Components}
\label{tab:stack}
\centering
\small
\begin{tabular}{p{0.28\columnwidth}p{0.62\columnwidth}}
\toprule
\textbf{Component} & \textbf{Implemented role}\\
\midrule
React 19, Vite & Browser-based workflow and production frontend build\\
Leaflet, Turf.js & Map display, polygon interaction, and geodesic area\\
FastAPI, Pydantic & HTTP API, request validation, and response schemas\\
SQLAlchemy, Alembic & ORM models and database schema migration\\
scikit-learn, Joblib & Bundled configuration and yield model inference\\
SHAP & Optional tree-model feature attributions\\
pmdarima & Seasonal rainfall forecast service\\
Open-Meteo & Historical precipitation retrieval\\
Groq (optional) & Structured natural-language advisory\\
\bottomrule
\end{tabular}
\end{table}

The frontend requires network access for rainfall retrieval, map imagery,
and the optional advisory. Backend settings provide the database URL and
JWT configuration; the advisory additionally requires a Groq API key.
The map and external rainfall calls transmit location coordinates to
their respective providers. The advisory request includes the society
name and submitted summary metrics. These integrations are operational
dependencies and should be considered in deployment privacy and
availability reviews.

\section{Evaluation and Results}
\label{sec:evaluation}

This paper reports implementation behavior, not a hydrological field
trial. No measurements of harvested runoff, storage operation, groundwater
response, or audited society savings were available for independent
validation. Model performance figures from saved offline experiments
are not reported as runtime evaluation results. Likewise, no forecast
accuracy, classifier generalization, SHAP stability, execution-latency,
or payback outcome is inferred from the mere availability of a model or
service.

The current backend test suite covers selected software invariants. The
tank-sizer tests exercise year-length handling, water-volume
reconciliation, zero-capacity behavior, monotonic capacity comparisons,
and finite recharge overflow. Monte Carlo integration tests exercise the
use of supplied monthly rainfall, invalid rainfall rejection, and the
effect of different rainfall inputs on downstream results. ROI tests
cover the deterministic calculation contract. These tests verify
specific code paths with controlled inputs; they do not establish the
accuracy of external rainfall data, the correctness of project
hydrogeological thresholds, real-world system reliability, or financial
returns. During this revision, all 17 backend \texttt{unittest} cases and
all 18 frontend unit tests passed; the frontend production build and
ESLint checks also completed successfully. The build emitted a
bundle-size warning, which does not change the calculation evidence
described above.

\begin{table}[htbp]
\caption{Evidence Scope for Implemented Outputs}
\label{tab:evidence}
\centering
\small
\begin{tabular}{p{0.28\columnwidth}p{0.60\columnwidth}}
\toprule
\textbf{Output} & \textbf{Evidence and interpretation}\\
\midrule
Demand and yield & Deterministic equations and schemas are implemented; site inputs and coefficients remain assumptions.\\
Rainfall forecast & API and seasonal model path are implemented; current holdout accuracy is not reported.\\
Configuration class & Bundled model is trained on synthetic rule-derived labels; external generalization is unverified.\\
Yield regression & Training target is physics-generated; independent observed-yield performance is unverified.\\
Tank sizing & Daily water-balance logic is covered by focused unit tests; no installed tank was measured.\\
ROI scenarios & Financial services are implemented and tested with controlled inputs; no realized savings are claimed.\\
\bottomrule
\end{tabular}
\end{table}

\section{Discussion and Limitations}
\label{sec:discussion}

WaterWallet's principal value is workflow integration: society-level
inputs, geospatial catchment estimation, rainfall data, site screening,
and financial scenarios are presented together. Its separation of
deterministic constraints from model-based recommendations makes the
decision path more inspectable than a single opaque score. The daily tank
simulator also distinguishes useful water, recharge overflow, lost
overflow, and unmet demand rather than treating all roof runoff as
financially recoverable.

Several limitations constrain interpretation:
\begin{enumerate}
  \item \textbf{Synthetic configuration labels.} The classifier learns
    rule-generated labels from a synthetic dataset. It cannot establish
    suitability for actual aquifers or discover site relationships
    absent from those rules. Independent surveys and expert-labelled
    sites are needed.
  \item \textbf{Unvalidated safety thresholds.} Coastal distance,
    permeability class, borewell presence, and borewell depth are
    simplified screening inputs. The gate's configured cutoffs are not a
    substitute for local hydrogeological investigation or regulatory
    review.
  \item \textbf{No measured yield data.} The yield regressor is trained
    against targets calculated from the same physics equation it
    approximates. No observed rooftop inflow dataset is used to show an
    advantage over that equation.
  \item \textbf{Rainfall provenance and uncertainty.} Historical
    retrieval may use a regional normal fallback; forecast-history
    retrieval may synthesize a regional series. Responses expose this
    provenance, but regional fallback data are not site-specific and
    the paper does not claim forecast accuracy for a given society.
  \item \textbf{Financial simplifications.} Deterministic ROI uses an
    annualized yield cap rather than daily storage. Monte Carlo sampling
    perturbs monthly rainfall independently and does not model
    multi-year drought sequences or recharge value. Cost, tariff,
    escalation, and maintenance inputs require local quotations and
    review.
  \item \textbf{Partial persistence and external services.} Ordinary
    society and billing flows use in-memory state despite the presence
    of database schemas and selected database endpoints. Rainfall,
    imagery, and optional advisory functions depend on external
    services.
  \item \textbf{Feasibility-screening boundary.} WaterWallet does not
    perform structural roof assessment, detailed filter design, water
    quality analysis, groundwater modelling, or regulatory approval.
\end{enumerate}

Future work should prioritize collecting real rainfall and harvested
volume observations from monitored sites; validating the screening rules
with qualified hydrogeologists; publishing a temporally held-out
forecast evaluation with data provenance; and comparing the yield model
against independently measured outcomes. The financial analysis could
be unified with daily tank balances and site-specific cost quotations,
while recharge benefits should remain excluded until defensible
hydrogeological measurements are available. Completing transactional
database persistence and retaining rainfall-source metadata with saved
analysis snapshots would also improve operational reliability.

\section{Conclusion}
\label{sec:conclusion}

WaterWallet is an implemented full-stack prototype for residential RWH
feasibility analysis in the Mumbai Metropolitan Region. It links demand
and billing inputs, map-derived catchment areas, rainfall retrieval and
forecasting, a safety-gated configuration classifier, a physics-informed
yield calculation, daily tank sizing, and deterministic and
probabilistic financial analyses. The architecture combines React and
FastAPI services with optional model explanations and an external
language-model advisory.

The current evidence supports describing the software components and
their tested calculation paths, but not claiming field performance or
generalizable predictive accuracy. Synthetic classifier labels,
equation-derived yield targets, non-site-specific rainfall fallbacks, simplified
financial assumptions, and incomplete persistence remain material
limitations. WaterWallet should therefore be used as a pre-consultation
decision aid, with final RWH configuration and construction determined
through site-specific professional assessment.

\begin{thebibliography}{00}

\bibitem{bis1172}
Bureau of Indian Standards, \textit{IS 1172:1993: Code of Basic
Requirements for Water Supply, Drainage and Sanitation}, New Delhi,
India, 1993.

\bibitem{openmeteo}
P. Zippenfenig, ``Open-Meteo.com Weather API,'' software, 2023,
doi: 10.5281/ZENODO.7970649. [Online]. Available:
\url{https://open-meteo.com/en/docs/historical-weather-api}

\bibitem{hyndman}
R. J. Hyndman and G. Athanasopoulos, \textit{Forecasting: Principles
and Practice}, 3rd ed. Melbourne, Australia: OTexts, 2021. [Online].
Available: \url{https://otexts.com/fpp3/}

\bibitem{lundberg}
S. M. Lundberg and S.-I. Lee, ``A unified approach to interpreting
model predictions,'' in \textit{Advances in Neural Information
Processing Systems}, vol. 30, 2017, pp. 4765--4774.

\end{thebibliography}

\end{document}
